\scanpage{181}\mspage{36}%
% archive sheet number: 181
% One continuous running paragraph; author number at top center, archive pencil number at lower left.
be able to be paraphrased in the setting of excellent schemes, with liberal use of étale tubular neighborhoods.
In both cases, the ``fine homotopy type'' of a tame space, resp. of an excellent scheme, would be defined by passage to the limit starting from the
% L7a: the opening quotation marks stand to the left of the raised prefix; its guide reaches the start of types.
``\inserted{$\unsure{\mathcal{P}|P}$-}homotopy types'' associated with finer and finer equisingular stratifications (with nonsingular strata) $\mathcal{P}$.
This \inserted{``fine''} homotopy type would embody knowledge, not of \unsure{locally|locally} constant sheaves or \inserted{$n$-}stacks alone, but
% L12a: complete interlinear parenthesis; the long guide descends before la.
\inserted{(via a passage to the \unsure{inductive} limit)} knowl%
% L12 nontext point at approximately (714,1189), below the ai region; origin uncertain.
\raisebox{0pt}[0pt][0pt]{\parbox{0pt}{\figure{\begin{tikzpicture}[overlay,x=1cm,y=1cm]\fill (-0.04,-0.28) circle (0.35pt);\end{tikzpicture}}}}%
edge of
% The visible rule begins within the o of tous, approximately (1171,1167), and ends near (1289,1164).
a\uline{ll}. And it would depend, in a suitable sense,
functorially%
% L14 nontext point at approximately (700,1366), below/right of fonctoriellement; origin uncertain.
\raisebox{0pt}[0pt][0pt]{\parbox{0pt}{\figure{\begin{tikzpicture}[overlay,x=1cm,y=1cm]\fill (0.16,-0.12) circle (0.3pt);\end{tikzpicture}}}}
on $X$. In%
% L14 nontext point at approximately (1170,1369), below the end of Dans; origin uncertain.
\raisebox{0pt}[0pt][0pt]{\parbox{0pt}{\figure{\begin{tikzpicture}[overlay,x=1cm,y=1cm]\fill (0,-0.27) circle (0.45pt);\end{tikzpicture}}}}
the case of a%
% L14 nontext point at approximately (1627,1352), below/right of d'un, near the right margin; origin uncertain.
\raisebox{0pt}[0pt][0pt]{\parbox{0pt}{\figure{\begin{tikzpicture}[overlay,x=1cm,y=1cm]\fill (0.45,-0.13) circle (0.25pt);\end{tikzpicture}}}}
scheme of finite type over an algebraically closed field say, the strongest cohomological or homotopical
% Literal ruled span: from within the é of théorèmes, approximately (671,1580), to (1328,1557), before cohomologique.
th\uline{eorems of finiteness} would be expressed precisely in terms of the fine homotopy types, and would say that the \unsure{ordinary|ordinary} homotopy types which are their constituent elements are essentially ``finite \unsure{polyhedra|polyhedra}'' -- and $\overline{m}$ compact manifolds with boundary -- or more precisely,
